\chapter{The evidence classes}
\label{ch:classes}

\section*{What this chapter is for}

Preparation moves some kinds of evidence and not others. The distinction
decides where a week is worth spending, and it is the reason this book opens
with a fit check rather than a study plan.

\section{Eight classes}

Everything a hiring process weighs falls into one of these. The third column is
the one that matters.

\begin{kittable}{@{}lll@{}}
\textbf{Class} & \textbf{What moves it} & \textbf{How long} \\
\midrule
Knowledge & study & days to weeks \\
Rehearsal & practice & days \\
Artefacts & building & weeks \\
Research & reading & hours \\
Consistency & discipline & hours \\
\midrule
Demonstrated scope & doing the job & months to years \\
Tenure pattern & staying & years \\
Structural fit & nothing & --- \\
\end{kittable}

The line across the middle is the whole chapter. Above it is what a preparation
kit can serve. Below it is what a preparation kit can only \emph{report}.

\section{The asymmetry, and what follows from it}

Effort above the line converts into evidence within the timescale of a hiring
process. Effort below it does not, however much of it there is.

\judgement{This produces a specific and very common failure. The classes above
the line are legible, they respond to effort, and working on them feels like
progress. So a candidate who is actually blocked below the line keeps preparing
\dash{} more study, another artefact, better answers \dash{} because that is the
part that moves when pushed. The preparation is real and the blocker is
elsewhere, and the loop can run several times before anybody says so.}

The remedy is cheap and it is an ordering: \textbf{check the bottom three
first}. Twenty minutes, before the week is committed.

\begin{kitmethod}
\textbf{Structural fit} \dash{} Chapter~\ref{ch:fit}. Six questions against the
posting.

\textbf{Tenure pattern} \dash{} look at your own history the way a stranger
will: the shape of the last few roles and their lengths, with no context
attached, because none will be attached when it is read.

\textbf{Demonstrated scope} \dash{} for the three pieces of work you will be
asked about, ask who decided each was the problem. Chapter~\ref{ch:staff} has
the test; Chapter~\ref{ch:deep-dive} has what to do with the answer.

Then, and only then, decide where the week goes.
\end{kitmethod}

\kitfig{evidence-classes}{% Chapter 13. The line is the content: above it, preparation converts into
% evidence inside a hiring process; below it, it does not.
\begin{tikzpicture}[node distance=3pt]
  \node[kitbox, minimum width=34mm] (know) {Knowledge\\{\scriptsize study \dash{} days}};
  \node[kitbox, minimum width=34mm, right=of know] (reh) {Rehearsal\\{\scriptsize practice \dash{} days}};
  \node[kitbox, minimum width=34mm, right=of reh] (art) {Artefacts\\{\scriptsize building \dash{} weeks}};
  \node[kitbox, minimum width=34mm, below=of know] (res) {Research\\{\scriptsize reading \dash{} hours}};
  \node[kitbox, minimum width=34mm, below=of reh] (con) {Consistency\\{\scriptsize discipline \dash{} hours}};
  \node[kitsoft, minimum width=34mm, below=of art] (gap) {\itshape a week of work\\\itshape moves all five};

  \coordinate (l) at ($(res.south west)+(0,-6mm)$);
  \draw[kitwarn, line width=1pt, dashed] ($(res.south west)+(-2mm,-6mm)$) -- ($(gap.south east)+(2mm,-6mm)$);
  \node[kitlabel, text=kitwarn, anchor=west] at ($(res.south west)+(-2mm,-3.2mm)$)
    {preparation reaches this far};

  \node[kitwarnbox, minimum width=34mm, below=13mm of res] (scope)
    {Demonstrated scope\\{\scriptsize doing the job \dash{} months}};
  \node[kitwarnbox, minimum width=34mm, below=13mm of con] (ten)
    {Tenure pattern\\{\scriptsize staying \dash{} years}};
  \node[kitwarnbox, minimum width=34mm, below=13mm of gap] (fit)
    {Structural fit\\{\scriptsize nothing moves it}};
\end{tikzpicture}
}{The eight classes. Above the dashed line, a week of work converts into evidence inside a hiring process. Below it, nothing you do this month reaches.}

\section{What to do with a class you cannot move}

Three options, and one of them is not on most people's list.

\textbf{Accept and prepare the answer.} A below-the-line weakness usually
surfaces as a question. Having the honest version prepared is worth doing, and
Chapter~\ref{ch:calibration} is about the difference between an objection that
an answer can settle and one that it cannot.

\textbf{Change the target.} A level down, a different segment, a company whose
constraint is different. This is not a retreat; it is the same work applied
where the immovable class is not in the way.

\textbf{Fix it, on a different timescale.} If the blocker is tenure or
demonstrated scope, the fix is months or years of doing something, and it is
usually where you already are. \judgement{Deciding to stay somewhere for two
years in order to change what your evidence says is a legitimate career move
and it is almost never described as one. It looks like inaction and it is
repair.}

\section{Where the classes appear in the loop}

\begin{kittable}{@{}lX@{}}
\textbf{Stage} & \textbf{What it weighs} \\
\midrule
Screen & consistency, structural fit \\
Technical & knowledge, rehearsal \\
Hiring manager & knowledge, research, artefacts \\
Build round & rehearsal, knowledge \\
Deep-dive & demonstrated scope \\
Behavioural & demonstrated scope, rehearsal \\
Executive & anything, including classes earlier stages already passed \\
\end{kittable}

The last row is the one to plan for. An executive round is not bound by what
earlier stages accepted, and a concern that was answered at one level can be
weighed again at another by somebody with different authority.
\judgement{Treating a raised concern as settled once it has been accepted is
the commonest error in estimating your own position, and it is an error in the
direction of optimism.}

\begin{drill}
Score yourself out of five on each of the eight classes for the role in front
of you, and write one line of evidence for each score. Then look only at the
bottom three.

If all three are strong, prepare with confidence: the process is winnable and
the work above the line decides it. If one is weak, you now know which question
is coming. If two are weak, the honest reading is that this particular role is
practice, and practice is a legitimate thing to buy as long as you know that
is what you are buying.
\end{drill}

\section*{What this chapter does not claim}

The eight classes are a framework, not a finding. No study measured them and
none is cited. They are offered because the distinction between evidence that
responds to a week of work and evidence that does not is real, and because
having the distinction changes how the week gets spent.

The stage-to-class table is an inference from what each stage is structurally
able to observe. It is a prior worth holding loosely.
